\chapter{Sizing and Selecting a Control Transformer}
\label{app:transformer-sizing}

\textbf{Applies to:} control power transformers (CPTs) that take plant power (for example 480~VAC) and make 120~VAC for panel control circuits --- power supplies, the PLC, contactor coils, fans, lights, and a service receptacle.

\section{The Short Version}
\begin{enumerate}[itemsep=2pt]
  \item Decide how the transformer connects (primary and secondary voltage).
  \item Add up the \textbf{sealed (steady) VA} of everything on 120~VAC.
  \item Add 25\% and round up to the next standard size.
  \item Check that the transformer can handle the \textbf{inrush} when coils pull in.
  \item Size the primary and secondary fuses.
  \item Pick a part and wire it per the nameplate.
\end{enumerate}

\section{Step-by-Step}
\subsection*{Step 1 --- Decide how it connects}
A 480~VAC \emph{delta} system has three hot lines and no neutral. A single-phase control transformer is connected \textbf{line-to-line} across two of the phases (for example L1 and L2), so its primary sees 480~V. The secondary is 120~V. One side of the secondary (X2) is bonded to ground to make the grounded (neutral) conductor; the other side (X1) is the fused hot.

Most control transformers have dual-voltage windings (for example 240$\times$480 primary, 120$\times$240 secondary). You set the voltage with jumpers shown on the nameplate.

\subsection*{Step 2 --- Add up the sealed VA}
``Sealed'' VA is what each device draws once it is running. Get it from the datasheet; for power supplies use input VA (output power $\div$ efficiency $\div$ power factor), not output watts. Only count loads that can be on at the same time.
\[
  \text{VA}_{\text{sealed}} = \sum \text{VA}_{\text{each device}}
\]

\subsection*{Step 3 --- Add margin and round up}
Add 25\% for growth and continuous loading, then pick the next standard size up. Common sizes are 50, 75, 100, 150, 250, 350, 500, 750, 1000, 1500, 2000, 3000, and 5000~VA.
\[
  \text{VA}_{\text{required}} = 1.25 \times \text{VA}_{\text{sealed}}
\]

\subsection*{Step 4 --- Check inrush}
When a contactor or relay coil first pulls in, it draws several times its sealed VA for a few cycles. If many coils pull in at once, the transformer's output voltage dips; if it dips too far, coils chatter or don't close. Combine the sealed load with the inrush of the coils that can energize at the same moment:
\[
  \text{VA}_{\text{inrush selection}} = \sqrt{\left(\text{VA}_{\text{sealed}}\right)^2 + \left(\text{VA}_{\text{inrush}}\right)^2}
\]
Then look up the chosen transformer in the manufacturer's \emph{regulation} (inrush) chart, usually in the 90\% secondary voltage column at 40\% power factor. The transformer's listed inrush rating must be at least your calculated value. If not, go up one size.

\subsection*{Step 5 --- Size the fuses}
The National Electrical Code (NEC 450.3(B), and 430.72(C) for motor control circuits) sets the maximum fuse size. For a transformer with \emph{both} primary and secondary fuses:
\begin{center}
\small
\begin{tabular}{@{}lll@{}}
\toprule
\rowcolor{tblhead}\textbf{Where} & \textbf{Rated current} & \textbf{Maximum fuse size} \\ \midrule
Primary & any (with secondary fuse) & 250\% of primary current \\
Secondary & 9 A or more & 125\% (next standard size allowed) \\
Secondary & less than 9 A & 167\% \\
Primary only (no secondary fuse) & 9 A or more & 125\% \\
Primary only (no secondary fuse) & 2 A to 9 A & 167\% \\
Primary only (no secondary fuse) & less than 2 A & 300\% (500\% allowed by 430.72(C) for control circuits) \\
\bottomrule
\end{tabular}
\end{center}
\[
  I_{\text{primary}} = \frac{\text{VA}}{V_{\text{primary}}}
  \qquad\qquad
  I_{\text{secondary}} = \frac{\text{VA}}{V_{\text{secondary}}}
\]
Use \textbf{time-delay} fuses (for example Class CC time-delay) on the primary so the transformer's own magnetizing inrush doesn't blow them. Both primary lines are fused because neither leg of a 480~V delta is grounded. On the secondary, fuse only X1; X2 is the grounded conductor.

\section{Worked Example: 480~VAC Delta to 120~VAC}
A panel is fed from a 3-phase, 480~VAC delta supply and needs 120~VAC control power. The 120~VAC loads are (example values --- use your actual datasheets):
\begin{center}
\small
\begin{tabular}{@{}lrr@{}}
\toprule
\rowcolor{tblhead}\textbf{Load} & \textbf{Sealed VA} & \textbf{Inrush VA} \\ \midrule
24~VDC power supply, 240~W output (92\% efficient, PF 0.95) & 275 & --- \\
Modicon BMXCPS3500 PLC power supply & 75 & --- \\
Enclosure fan & 40 & --- \\
LED panel light & 15 & --- \\
Service receptacle (NEC counts 180~VA) & 180 & --- \\
Three contactor coils (10~VA sealed / 100~VA inrush each) & 30 & 300 \\ \midrule
\textbf{Total} & \textbf{615} & \textbf{300} \\
\bottomrule
\end{tabular}
\end{center}

\begin{enumerate}[itemsep=4pt]
  \item \textbf{Connection.} Single-phase transformer, primary across L1--L2 at 480~V, secondary 120~V with X2 bonded to ground.
  \item \textbf{Sealed VA} $= 615$~VA.
  \item \textbf{Margin:} $1.25 \times 615 = 769$~VA $\Rightarrow$ next standard size is \textbf{1000~VA}.
  \item \textbf{Inrush check} (all three contactors pulling in together):
  \[
    \sqrt{615^2 + 300^2} = \sqrt{378{,}225 + 90{,}000} \approx 684~\text{VA}
  \]
  Look up the 1000~VA model in the manufacturer's regulation chart (90\% column). Its inrush rating must be at least 684~VA; 1000~VA control transformers are normally rated well above that, so 1000~VA is fine. Confirm on the chart for the exact part.
  \item \textbf{Fuses.}
  \begin{itemize}[itemsep=1pt]
    \item Primary current: $1000 / 480 = 2.08$~A. Maximum with secondary protection: $2.5 \times 2.08 = 5.2$~A $\Rightarrow$ \textbf{5~A time-delay Class CC}, one in each primary line (2 fuses).
    \item Secondary current: $1000 / 120 = 8.33$~A. Less than 9~A, so the maximum is $1.67 \times 8.33 = 13.9$~A $\Rightarrow$ \textbf{10~A} fuse on X1 (leaves margin above the 8.33~A rating).
  \end{itemize}
  \item \textbf{Select a part.} For example, the \href{https://www.lamonde.com/ph1000mqmj.html}{Hammond Power Solutions PH1000MQMJ}: 1000~VA, 240$\times$480~V primary, 120$\times$240~V secondary. Jumper the primary for 480~V and the secondary for 120~V exactly as shown on its nameplate diagram. Add finger-safe fuse blocks for the primary and secondary fuses, or choose the factory fuse-kit version (PH1000MQMJ-FK) if available.
\end{enumerate}

\textbf{Result:} 1000~VA, 480~V to 120~V control transformer; 2 $\times$ 5~A time-delay primary fuses; 1 $\times$ 10~A secondary fuse on X1; X2 bonded to ground.

\section{Further Reading}
\begin{itemize}[itemsep=2pt]
  \item \href{https://www.elliottelectric.com/StaticPages/ElectricalReferences/Guides/control-transformer-selection-guide.aspx}{Elliott Electric --- Control transformer selection guide (inrush VA equation and regulation chart)}
  \item \href{https://literature.rockwellautomation.com/idc/groups/literature/documents/wp/1497-wp001_-en-p.pdf}{Rockwell Automation --- Proper selection of control circuit transformers (white paper)}
  \item \href{https://electrical-engineering-portal.com/overcurrent-protection-transformer-nec-450-3}{Electrical Engineering Portal --- Overcurrent protection of transformers (NEC 450.3)}
\end{itemize}
